% Slide 15 — Evaluation Metrics
\begin{frame}[shrink=2]{Evaluation Metrics}
  \begin{columns}[T,onlytextwidth]
    \column{0.56\textwidth}
    \begin{block}{What we measure (macro-averaged)}
      \small
      \begin{tabular}{@{}p{0.22\linewidth}p{0.72\linewidth}@{}}
        \toprule
        \textbf{Metric} & \textbf{Why it matters} \\
        \midrule
        Accuracy & overall correctness \\
        Precision & avoid false positives (per class) \\
        Recall & avoid missed tumors (per class) \\
        F1-macro & balances precision/recall equally across classes \\
        $\kappa$ & agreement beyond chance (robustness) \\
        \bottomrule
      \end{tabular}
    \end{block}

    \begin{block}{Ranking criterion}
      \small We primarily rank configurations by \keyresult{macro-F1} to ensure fair
      performance across the three tumor classes.
    \end{block}

    \column{0.42\textwidth}
    \centering
    \begin{tikzpicture}[
      node distance=0.5cm,
      box/.style={rectangle, rounded corners=2pt, draw=accent, fill=lightbg,
                  minimum height=0.7cm, minimum width=0.92\linewidth, align=center, font=\small},
      arr/.style={-{Stealth[length=2mm]}, thick, accent}
    ]
      \node[box] (train) {Train split (80\%)}; 
      \node[box, below=of train] (cv) {5-fold CV + tuning};
      \node[box, below=of cv] (test) {Held-out test (20\%)};
      \node[box, below=of test] (stats) {CI + significance tests};
      \draw[arr] (train) -- (cv);
      \draw[arr] (cv) -- (test);
      \draw[arr] (test) -- (stats);
    \end{tikzpicture}

    \vspace{0.3em}
    {\scriptsize Bootstrap CI, McNemar (paired test errors), Friedman (multi-model comparison).}
  \end{columns}
\end{frame}

\note{
Macro averaging treats all tumor classes equally --- important for clinical fairness.
F1-macro is the primary ranking metric.
}
